\section{Related work}
\label{sec:related-work}

We organize related work along several axes:\ the textbook tradition
the library mechanizes, prior efforts to mechanize fragments of
game theory in interactive theorem provers, current Lean projects in
the same neighborhood, and adjacent formalizations in dependent type
theory or higher-order logic that could plausibly serve as alternative
platforms.

\subsection{Textbook game theory}
\label{sec:rw:textbooks}

The library's substantive content is drawn from the standard
textbook canon.  \citet{vNM1944} laid the basis for the
expected-utility decision theory the library uses without
modification; the minimax theorem first appears
in~\citet{vN1928}.  \citet{Nash1950} and~\citet{Nash1951} are the
sources for the existence theorem in the form proved in
\Cref{thm:nash-mixed}.  \citet{Aumann1974, Aumann1987} introduce
correlated equilibrium and its EU equivalence with
incentive-compatibility under a correlation device.
\citet{Selten1965} introduces subgame perfection and
\citet{Selten1975} trembling-hand perfection, both of which the
library formalizes.  The one-shot deviation principle---the
standard characterization of subgame-perfect equilibrium in
finite perfect-information games---is the form proved in
\Cref{thm:odp}.
\citet{Harsanyi1967} formalizes Bayesian games and
type-conditional strategies.  \citet{Kuhn1953} is the source of
the behavioral-vs-mixed equivalence, originally proved under
perfect recall;~\citet{Aumann1964} extended it to infinite
extensive games.
\citet{Friedman1971} and~\citet{FudenbergMaskin1986} are the
classical references for the discounted repeated-game folk theorem
formalized in \Cref{thm:folk-discounted}.

For comprehensive surveys we draw
on~\citet{OsborneRubinstein1994},
\citet{MaschlerSolanZamir2013},~\citet{ShohamLeytonBrown2008},
and~\citet{Myerson1991}.  The mechanism-design layer follows
Myerson's revelation-principle exposition;
the auction layer follows~\citet{Krishna2010}.  Algorithmic
game theory, including the price of anarchy, is treated
in~\citet{NisanRoughgardenTardosVazirani2007}, with Roughgarden's
smoothness framework~\citep{Roughgarden2015}, his selfish-routing
analysis of the price of anarchy~\citep{Roughgarden2005}, and
Koutsoupias–Papadimitriou's original
PoA~\citep{Koutsoupias1999} all relevant.

The factored-observation stochastic-game
formalism~\citep{Kovarik2022} unifies extensive-form games,
partially observable Markov games, and
multi-agent reinforcement-learning settings; the library uses it
as the most general native presentation, while keeping EFG and
NFG as distinct languages because the textbook tradition treats
them so.  The sequence-form representation
of~\citet{KollerMegiddoVonStengel1994} is conspicuously absent
from the library; we discuss this in \Cref{sec:future}.

Multi-agent influence diagrams are due
to~\citet{KollerMilch2003} and have been further developed
in~\citet{PfefferGal2007} and~\citet{HammondEtAl2023}.  The
order-independence and frontier-commutation lemmas for MAID
compilation are partly formalized in our library
(\TT{MAID/FrontierLemmas.lean}); strategic-relevance and
d-separation analyses remain extension targets.

\subsection{Prior formalizations of game theory}
\label{sec:rw:formalization}

\paragraph{Coq.}\ \citet{Vestergaard2006} gives a Coq-formalized
constructive proof of sequential Nash equilibria.
\citet{LeRouxMartinDorelSmaus2017} formalize a Nash-equilibrium
existence theorem in both Coq and Isabelle, taking a constructive
route distinct from our Brouwer-based proof.
\citet{Bertot2016} formalized the Stable Marriage / Gale–Shapley
algorithm in Coq.

\paragraph{Isabelle/HOL.}\ Lange, Caminati, Kerber, Rowat, and
collaborators have produced a substantial body of mechanized
auction theory in Isabelle, including VCG, combinatorial
auctions, and revenue-equivalence
fragments~\citep{KerberLangeRowat2016, CaminatiKerberLangeRowat2014}.
Their setting is more directly continuous than ours, since
Isabelle's analysis library makes real-valued bid spaces
ergonomic.  Isabelle's CryptHOL~\citep{Lochbihler2016} formalizes
discrete probability and probabilistic programs and has been
used for cryptographic-protocol verification rather than
game-theoretic equilibria.

\paragraph{Lean.}\ The discrete-probability monad \TT{PMF}, the
standard simplex and its convex/topological API, finite-type
infrastructure, and the surrounding analysis the library depends
on are all in mathlib~\citep{mathlib2020}; the Brouwer fixed-point
theorem itself is supplied by the standalone
\TT{fixed-point-theorems-lean4}
package~\citep{harfeFixedPointTheorems}.

The closest neighboring project is
\TT{EconCSLib}~\citep{EconCSLib2026}, a public Lean~4 library and
knowledge base for computational economics.  Its current scope
includes strategic, extensive, and coalitional games; social choice
and fair division; matching; auctions and mechanism design; and
supporting fixed-point, minimax, and linear-programming material.
The overlap with the present library is therefore substantial:
equilibrium concepts, zero-sum games, mechanism design, voting,
auctions, matching, and cooperative-game primitives all appear in
both developments.

The projects nevertheless have different design centers.
\TT{EconCSLib} is organized as a broad computational-economics
library, with lightweight domain-specific APIs, locally stated
assumptions, and a large blueprint/knowledge-base layer.  The
present library is organized around semantic interoperability:
\TT{GameForm} and \TT{KernelGame} are the common stochastic targets,
and the language layer proves compilation and bridge theorems from
NFG, EFG, MAID, FOSG, multi-round, and intrinsic presentations into
that target.  The two code bases are not API-compatible, but the
mathematical overlap is large enough that a shared
crosswalk of definitions, theorem statements, and intended adapters
would be valuable.  Such coordination would reduce duplicated proof
engineering and clarify which results should be shared, ported, or
kept deliberately separate because they serve different interfaces.

Several Lean developments focus on the Nash-existence slice.
\TT{math-xmum/Brouwer}~\citep{MathXMUMBrouwer2026} and the
\TT{Solo-ary/Game-Theory-Formalization}
artifact~\citep{SoloAryGameTheoryFormalization2026} both formalize
a Scarf/Brouwer route to mixed Nash equilibria in finite games.
These are directly relevant to our \Cref{thm:nash-mixed} proof
route, but their stated endpoint is narrower:\ fixed-point
infrastructure and Nash existence, not a multi-language semantic
library with theorem transport across representations.
\TT{MixedMatched/formalizing-game-theory}
~\citep{MixedMatchedFormalizingGameTheory} is an exploratory Lean
formalization of game-theoretic structures, including heterogeneous
strategy types and pure/mixed strategy profiles.  It is useful as a
modeling comparison, but it does not try to provide the language
layer, representation bridges, or theorem suite developed here.

Other Lean projects overlap with specific subfields rather than with
the whole library.  \TT{lean-social-choice}
~\citep{LeanSocialChoice} formalizes social-choice results including
Arrow's theorem; \TT{apportionmentlib}~\citep{Apportionmentlib2026}
formalizes apportionment theory, including Balinski--Young-style
impossibility results.  The combinatorial-games library
\citep{LeanCombinatorialGames} develops Conway-style terminating
two-player perfect-information games, nimbers, and surreal-number
connections.  It is deep game theory, but it is a different
formalism from the stochastic, expected-utility, partial-information
games treated here.

Among the Lean projects we found, the feature combination emphasized
in this paper remains distinctive:\ a protocol/preference split, a
multi-language layer, FOSG and MAID/EFG bridges, the Kuhn theorem at
an observation-model level, and mechanism-design constructions all
inside one semantic framework.  This is the comparison that matters
for the claims made here; it is not a claim that other Lean
formalizations are merely isolated or that future coordination is
unnecessary.

\paragraph{Other systems.}\ ACL2, HOL4, and Mizar each have
isolated game-theoretic results we are aware of but no substantial
library; we omit them for space.

\subsection{Verified probability and probabilistic programming}
\label{sec:rw:probability}

The library's reliance on \TT{PMF} places it in a tradition of
discrete probability formalizations.  \citet{HurdMcIverMorgan2005}
mechanized probabilistic guarded commands in HOL, and
\citet{Audebaud2009} formalized proofs of randomized algorithms
in Coq, both with semantic foundations that anticipate the
discrete Giry monad.  \citet{Tassarotti2019} developed a
separation logic for concurrent randomized programs in Coq.
\citet{Lochbihler2016} (CryptHOL) does similar
work in Isabelle, with cryptographic applications.  Mathlib's
\TT{PMF}~\citep{mathlib2020} provides the API on which our
library builds.

The measure-theoretic side is more developed in mathlib than the
discrete side, with substantial integration across the Lean Borel,
Lebesgue, and product-measure infrastructures; an extension of the
present library over mathlib's measure theory would draw on this.

\subsection{Mechanized programming-language semantics}
\label{sec:rw:pl}

The library's language layer borrows several patterns from
mechanized programming-language semantics:\ structural
operational semantics, denotational compilation into a common
target, simulation/refinement~\citep{Milner1989, AbramskyJung1994,
HarperPLFA, BentonHurKennedyMcBride2014}, and CompCert-style
verified-translation pipelines~\citep{Leroy2009}.  The
``language as syntax + compiler to a common semantics'' framing
is direct PL practice; what is unusual in our setting is that the
common semantics is a strategic object (a kernel game) rather
than a behavioral one.

\subsection{Categorical game theory}
\label{sec:rw:category}

A growing body of work treats games and strategies as objects in
a category with a compositional algebra of game forms.
\citet{Pavlovic2009, Hyland2010} are early examples;
\citet{GhaniHedgesWinschelZahn2018} introduce \emph{open games}
and the bidirectional ``selection function'' approach;
\citet{Joyal1977} pioneered the categorical view of two-player
games (in a logical setting).  Our protocol-map and simulation
infrastructure is inspired by this tradition without committing
to any of its specific frameworks.  An open-games-as-coalgebras
approach would be a natural reformulation; we did not take it
because the kernel-game record is already small enough that
adding a categorical superstructure would, on net, cost
ergonomics.

\subsection{Mechanized formal verification of equilibria in protocols}
\label{sec:rw:protocols}

A different line of work mechanizes
equilibrium reasoning about cryptographic and economic protocols
directly.  \citet{HalpernPass2008, AbrahamDolevGonenHalpern2006}
and successors formalize rationality-based protocol analysis and
incentive-compatible distributed algorithms.

Recent Lean case studies are in the same spirit.  The
\TT{VerifiedHandshake} artifact~\citep{VerifiedHandshake2026}
mechanizes a commitment wrapper for Stag-Hunt-ordered games and
audits the axiom footprint of the proved statements.  It is not a
general game-theory library; it is a sharp example of a
protocol-specific equilibrium theorem.  Lean AMM work is similarly
specific:\ \citet{PuscedduBartoletti2024AMM} and the associated
\TT{lean4-amm} repository~\citep{Lean4AMM} formalize automated
market-maker state and economic properties, while
\TT{amm-axioms-lean}~\citep{AMMAxiomsLean2026} derives
constant-product and weighted-geometric-mean AMM behavior from swap
axioms.  These projects should be viewed as downstream or adjacent
economic-protocol formalizations, not as substitutes for a reusable
library of game representations.  They typically commit to one
domain model and prove protocol-specific results.  Our library is
meant to make that style of work less bespoke by supplying the
strategic-form infrastructure, equilibrium predicates, stochastic
semantics, and representation bridges that such developments would
otherwise build by hand.
